\documentclass[12pt, a4paper]{article}
\usepackage[utf8]{inputenc}
\usepackage[T1]{fontenc}
\usepackage[romanian]{babel}
\usepackage{amsmath}
\usepackage{amssymb}
\usepackage{geometry}
\geometry{a4paper, margin=2cm}
\usepackage{listings}
\usepackage{xcolor}

\lstset{
    language=Python,
    basicstyle=\ttfamily\small,
    keywordstyle=\color{blue},
    stringstyle=\color{red},
    commentstyle=\color{green!50!black},
    showstringspaces=false,
    breaklines=true,
    frame=single,
    numbers=left,
    numberstyle=\tiny\color{gray}
}


\title{Rezolvare Teme - Algoritmi de Primalitate \c{s}i Factorizare}
\author{}
\date{}

\begin{document}

\maketitle

\section*{Tema 1: Implementarea algoritmului Solovay-Strassen}

Algoritmul Solovay-Strassen este un test de primalitate probabilist care verific\u{a} dac\u{a} un num\u{a}r $n$ satisface condi\c{t}ia de corela\c{t}ie \^{i}ntre simbolul Jacobi $\left(\frac{b}{n}\right)$ \c{s}i congruen\c{t}a de tip Euler $b^{\frac{n-1}{2}} \pmod n$.

\begin{lstlisting}
import random

def simbol_jacobi(a, n):
    if a == 0: return 0
    if a == 1: return 1
    res = 1
    if a < 0:
        a = -a
        if n % 4 == 3: res = -res
    while a != 0:
        while a % 2 == 0:
            a //= 2
            if n % 8 in [3, 5]: res = -res
        a, n = n, a
        if a % 4 == 3 and n % 4 == 3: res = -res
        a %= n
    return res if n == 1 else 0

def test_Solovay_Strassen(n, nr_incercari):
    if n == 2: return True
    if n < 2 or n % 2 == 0: return False
    
    for _ in range(nr_incercari):
        b = random.randint(2, n - 1)
        if (n % b == 0): return False 
        
        jacobi = simbol_jacobi(b, n) % n
        putere = pow(b, (n - 1) // 2, n)
        
        if putere != jacobi:
            return False
    return True
\end{lstlisting}

\end{document}